% GENERATED FILE. Edit canonical lessons, then run npm run generate:books.
% canonical-source-hash: fnv1a64:ec79be8e3baa504c
% canonical-lessons: PT-C336-tema, PT-C336-questao, PT-C336-debate, PT-C336-debater, PT-C336-argumento

\chapter{O debate}
\label{ch:pt-o-debate}
\hlchaptermodality{336}

\begin{chapteropening}
\textbf{By the end of this chapter:} \emph{I can take part in a debate: name the topic, raise an issue and put forward an argument.}

Complete the last lesson of chapter 336: I can take part in a debate: name the topic, raise an issue and put forward an argument.
\end{chapteropening}

\section[o tema]{o tema — the topic, the subject; the theme}

\label{lesson:PT-C336-tema}



\begin{quote}
Before the new one: say the Portuguese for to give an opinion, then the Portuguese for to maintain.
\end{quote}

\subsection*{You'll want to know: o tema}
\textbf{o tema} — \textquotedblleft{}the topic, the subject; the theme\textquotedblright{}.

\begin{quote}
Qual é o tema da reunião de hoje?
\end{quote}

\begin{quote}
\emph{What is the topic of today's meeting?}
\end{quote}

\emph{Mudar de tema} — to change the subject.

\begin{cousinweb}
\emph{tema} is a learned borrowing from Latin \emph{thema} (EtymDB), which took it from Greek \emph{théma} \textquotedblleft{}something set down, a proposition\textquotedblright{}, built on \emph{tithénai} \textquotedblleft{}to put\textquotedblright{}.
\end{cousinweb}

\begin{grammarlens}[title={What you've built}]
One more word for telling, explaining and describing at B1.
\end{grammarlens}

\subsection*{Your turn}
\begin{itemize}
\raggedright
  \item \emph{Say it:} \emph{o tema}
  \item \emph{Say it:} \emph{o tema}, once more
  \item \emph{Say it:} a sentence of your own with \emph{o tema}
  \item \emph{Recall:} say the Portuguese for to give an opinion, then the Portuguese for to maintain, then say \emph{o tema} again
\end{itemize}

\begin{tcolorbox}[breakable,colback=teal!4,colframe=teal!35!black,title={Before you move on}]
What does o tema mean? (\textquotedblleft{}The topic, the subject; the theme\textquotedblright{}.) Say it once more.
\end{tcolorbox}
\section[a questão]{a questão — the issue, the matter; the question}

\label{lesson:PT-C336-questao}



\begin{quote}
Before the new one: say the Portuguese for to maintain, then the Portuguese for the topic.
\end{quote}

\subsection*{You'll want to know: a questão}
\textbf{a questão} — \textquotedblleft{}the issue, the matter; the question\textquotedblright{}.

\begin{quote}
A questão não é o dinheiro, é o tempo.
\end{quote}

\begin{quote}
\emph{The issue is not the money, it is the time.}
\end{quote}

\emph{Uma pergunta} is something you ask; \emph{uma questão} is an issue to discuss. \emph{Em questão} — in question.

\begin{cousinweb}
\emph{questão} is a learned borrowing from Latin \emph{quaestiō} \textquotedblleft{}a seeking, an inquiry\textquotedblright{} (EtymDB, Lewis \& Short), built on \emph{quaerere} \textquotedblleft{}to seek, to ask\textquotedblright{}.
\end{cousinweb}

\begin{grammarlens}[title={What you've built}]
One more word for telling, explaining and describing at B1.
\end{grammarlens}

\subsection*{Your turn}
\begin{itemize}
\raggedright
  \item \emph{Say it:} \emph{a questão}
  \item \emph{Say it:} \emph{a questão}, once more
  \item \emph{Say it:} a sentence of your own with \emph{a questão}
  \item \emph{Recall:} say the Portuguese for to maintain, then the Portuguese for the topic, then say \emph{a questão} again
\end{itemize}

\begin{tcolorbox}[breakable,colback=teal!4,colframe=teal!35!black,title={Before you move on}]
What does a questão mean? (\textquotedblleft{}The issue, the matter; the question\textquotedblright{}.) Say it once more.
\end{tcolorbox}
\section[o debate]{o debate — the debate}

\label{lesson:PT-C336-debate}



\begin{quote}
Before the new one: say the Portuguese for the topic, then the Portuguese for the issue.
\end{quote}

\subsection*{You'll want to know: o debate}
\textbf{o debate} — \textquotedblleft{}the debate\textquotedblright{}.

\begin{quote}
Ontem houve um debate na televisão sobre esta questão.
\end{quote}

\begin{quote}
\emph{Yesterday there was a debate on television about this issue.}
\end{quote}

\emph{Um debate aceso} — a heated debate.

\begin{cousinweb}
\emph{debate} is built on the verb \emph{debater}, which is usually traced to Old French \emph{debatre} \textquotedblleft{}to fight, to contend, to debate\textquotedblright{} (EtymDB lists the French verb), a formation on Latin \emph{battuere} \textquotedblleft{}to beat\textquotedblright{} (Lewis \& Short): ideas are knocked back and forth.
\end{cousinweb}

\begin{grammarlens}[title={What you've built}]
One more word for telling, explaining and describing at B1.
\end{grammarlens}

\subsection*{Your turn}
\begin{itemize}
\raggedright
  \item \emph{Say it:} \emph{o debate}
  \item \emph{Say it:} \emph{o debate}, once more
  \item \emph{Say it:} a sentence of your own with \emph{o debate}
  \item \emph{Recall:} say the Portuguese for the topic, then the Portuguese for the issue, then say \emph{o debate} again
\end{itemize}

\begin{tcolorbox}[breakable,colback=teal!4,colframe=teal!35!black,title={Before you move on}]
What does o debate mean? (\textquotedblleft{}The debate\textquotedblright{}.) Say it once more.
\end{tcolorbox}
\section[debater]{debater — to debate, to discuss}

\label{lesson:PT-C336-debater}



\begin{quote}
Before the new one: say the Portuguese for the issue, then the Portuguese for the debate.
\end{quote}

\subsection*{You'll want to know: debater}
\textbf{debater} — \textquotedblleft{}to debate, to discuss\textquotedblright{}.

\begin{quote}
Os alunos debateram o tema durante uma hora.
\end{quote}

\begin{quote}
\emph{The students debated the topic for an hour.}
\end{quote}

\emph{Debater um tema} — to debate a topic. More formal than \emph{discutir}.

\begin{cousinweb}
\emph{debater} is usually traced to Old French \emph{debatre} \textquotedblleft{}to fight, to contend, to debate\textquotedblright{} (EtymDB), built from \emph{de-} and \emph{batre} \textquotedblleft{}to beat\textquotedblright{}, from Latin \emph{battuere} \textquotedblleft{}to beat, to strike\textquotedblright{} (Lewis \& Short).
\end{cousinweb}

\begin{grammarlens}[title={What you've built}]
One more word for telling, explaining and describing at B1.
\end{grammarlens}

\subsection*{Your turn}
\begin{itemize}
\raggedright
  \item \emph{Say it:} \emph{debater}
  \item \emph{Say it:} \emph{debater}, once more
  \item \emph{Say it:} a sentence of your own with \emph{debater}
  \item \emph{Recall:} say the Portuguese for the issue, then the Portuguese for the debate, then say \emph{debater} again
\end{itemize}

\begin{tcolorbox}[breakable,colback=teal!4,colframe=teal!35!black,title={Before you move on}]
What does debater mean? (\textquotedblleft{}To debate, to discuss\textquotedblright{}.) Say it once more.
\end{tcolorbox}
\section[o argumento]{o argumento — the argument (a reason given); the plot}

\label{lesson:PT-C336-argumento}



\begin{quote}
Before the new one: say the Portuguese for the debate, then the Portuguese for to debate.
\end{quote}

\subsection*{You'll want to know: o argumento}
\textbf{o argumento} — \textquotedblleft{}the argument (a reason given); the plot\textquotedblright{}.

\begin{quote}
É um bom argumento, mas não me convence.
\end{quote}

\begin{quote}
\emph{It is a good argument, but it does not convince me.}
\end{quote}

An \emph{argumento} is a reason you give; a quarrel is \emph{uma discussão}. It is also the plot or script of a film.

\begin{cousinweb}
\emph{argumento} is a learned borrowing from Latin \emph{argūmentum} \textquotedblleft{}proof, evidence\textquotedblright{} (EtymDB), built on \emph{arguere} \textquotedblleft{}to make clear, to prove\textquotedblright{} (Lewis \& Short).
\end{cousinweb}

\begin{grammarlens}[title={What you've built}]
One more word for telling, explaining and describing at B1.
\end{grammarlens}

\subsection*{Your turn}
\begin{itemize}
\raggedright
  \item \emph{Say it:} \emph{o argumento}
  \item \emph{Say it:} \emph{o argumento}, once more
  \item \emph{Say it:} a sentence of your own with \emph{o argumento}
  \item \emph{Recall:} say the Portuguese for the debate, then the Portuguese for to debate, then say \emph{o argumento} again
\end{itemize}

\begin{tcolorbox}[breakable,colback=teal!4,colframe=teal!35!black,title={Before you move on}]
What does o argumento mean? (\textquotedblleft{}The argument (a reason given); the plot\textquotedblright{}.) Say it once more.
\end{tcolorbox}
